% Teorema frontal aditivo. No sustituye ningún capítulo ni altera el censo 102.
% BEGIN HMT-MD-FRONTAL-MAXIMAL-SIGNATURES-20260904
\chapter*{Théorème générateur de l'holographie modulaire triadique}
\addcontentsline{toc}{chapter}{Théorème générateur de l'holographie modulaire triadique}
\markboth{Théorème générateur de l'holographie modulaire triadique}%
{Théorème générateur de l'holographie modulaire triadique}

La démonstration commence dans le domaine fini
\(\mathcal D_9=\{1,\ldots,9\}\). Deux curseurs orientés parcourent,
respectivement, les feuillets additif et multiplicatif. La division euclidienne
conserve résidu et quotient, et le catalogue APP produit
\[
 \sum_{i,j\in\mathcal D_9}R^+_{ij}=405,
 \qquad
 \sum_{i,j\in\mathcal D_9}R^\times_{ij}=459,
 \qquad 459-405=54.
\]
Le transport coordonné
\(c(a)=a\bmod9:\mathcal D_9\to\{0,\ldots,8\}\)
change la carte et conserve la classe. La composition effective d'APP, de TRIT et de
TPK conserve à chaque pas feuillet, régime, résidu, quotient, retenue,
direction, bord et mémoire. Le catalogue exhaustif établit la chaîne
finie initiale et apporte son certificat exécutable :
\begin{equation}
\boxed{
104\,976\ \text{paires de germes}
\xrightarrow{\;54\ \text{pas}\;}
468\,U_6
\xrightarrow{\;\rho_3\;}
243\,w_6\subset\mathbb F_3^6,
\qquad |\mathbb F_3^6|=729.}
\label{eq:frontal-censo-semillas}
\end{equation}
L'identité des fibres
\[
432\cdot144+18\cdot432+18\cdot1944=104\,976
\]
prouve l'exhaustivité de la première application. Les nombres \(468\), \(243\) et
\(729\) comptent les émissions, les mots réalisés et l'espace ambiant ternaire,
respectivement ; leurs types demeurent séparés par les flèches de
\eqref{eq:frontal-censo-semillas}.

\begin{theorem}[Chaîne génératrice unique]
Soit \(\mathfrak H_9\) la structure formée par APP, TRIT, TPK et sa fibre
enrichie. APP exécute les feuillets somme et produit et conserve la division
euclidienne
\[
i+j=R^+_{ij}+9q^+_{ij},
\qquad
ij=R^\times_{ij}+9q^\times_{ij}.
\]
TRIT oriente chaque transition dans les régimes elliptique, parabolique et
hyperbolique. TPK compose
\[
    \boxed{\mathcal U_t
    =\operatorname{Upd}_t\circ\operatorname{Tra}_t
     \circ\operatorname{Sel}_t:
     X_{\rm TPK}^{\rm enr}\longrightarrow X_{\rm TPK}^{\rm enr}.}
\]
\(\operatorname{Sel}_t\) relève le sélecteur tritique,
\(\operatorname{Tra}_t\) relève le transport cardinal et
\(\operatorname{Upd}_t\) met à jour résidu, quotient, retenue, signature,
cylindre, bord, registre et mémoire sans effacer la provenance.
L'holonomie nonadique prolonge les mots par
\[
w_6\longrightarrow w_{12}\longrightarrow w_{18}\longrightarrow w_{24}
\longrightarrow w_{30}\longrightarrow R_{36}\longrightarrow G_9,
\]
et satisfait
\[
    \Gamma_{9,t}=\mathcal U_{t+8}\circ\cdots\circ\mathcal U_t,\qquad
    g(t+9)=g(t),
    \qquad q_{\rm mem}(\Gamma_{9,t}\widehat x_t)
    =q_{\rm mem}(\widehat x_t)+1.
\]
En général, \(\Gamma_{9,t}\widehat x_t\ne\widehat x_t\) : le retour conserve
la phase et augmente la mémoire. L'ordre opératoriel est
\(\mathcal U_t\to\Gamma_9\to K_{\rm ph}\) ; \(K_{\rm ph}\) est une publication
réduite ultérieure, non le générateur. Les troncatures
compatibles forment un système inverse non vide et définissent un état
\(x_\infty\) de profondeur arbitraire.

Cet état engendre une structure discrète unique du continu et deux
publications naturelles :
\begin{equation}
x_\infty
\xrightarrow{\ (\mathcal P_{\mathrm{arq}},Q_\infty)\ }
\mathbb R\times(\mathscr C_W^{\mathrm{cal}})^{\mathbb N}.
\label{eq:frontal-doble-proyeccion}
\end{equation}
La première contracte des cylindres compatibles jusqu'à une valeur réelle ; la seconde
conserve l'incidence et la généalogie jusqu'aux constructions de Paley, de Witt,
de Golay, de Leech, de \(V^\natural\) et du groupe Monstre. Sur la même source sont
publiés les quatre caractères corrélés
\[
(\pi_{\mathrm{HMT}},\varphi_{\mathrm{HMT}},
e_{\mathrm{HMT}},\alpha_{\mathrm{HMT}}),
\]
et, au moyen de transductions typées, les lois d'action, de matière, de masse,
de gravité et d'information de la mécanique dimensionnelle.
\end{theorem}

\begin{proposition}[Reconstruction intrinsèque et autosuffisance probatoire]
\label{prop:frontal-autosuficiencia-probatoria}
Chaque résultat \(y\) publié dans ce volume est reconstruit au moyen du
reçu mathématique local
\[
 \mathscr R_y=
 \bigl(D_y,A_y,\tau_y,U_y,C_y,X_y^{\rm enr},
       y_{\rm HMT},\operatorname{Rec}_y,\mathcal F_y\bigr),
\]
où \(D_y\) est le domaine ; \(A_y\) spécifie les feuillets et les opérations
d'APP ; \(\tau_y\) le régime et l'orientation TRIT ; \(U_y\) la composition TPK
avec domaine, codomaine et action ; \(C_y\) les données conservées de résidu,
quotient, retenue, feuillet, voie, bord, mémoire et survie ;
\(X_y^{\rm enr}\) l'état enrichi de résidence ; \(y_{\rm HMT}\) la
sortie antérieure à toute identification externe ; \(\operatorname{Rec}_y\) la
reconnaissance conventionnelle ultérieure ; et \(\mathcal F_y\) un réfutateur de la
flèche substantielle. La suite
\[
 \mathcal D_9\longrightarrow\Omega_{\rm seed}
 \longrightarrow\mathcal U_6^{\rm cat}
 \longrightarrow\mathcal W_6
 \longrightarrow X_\infty^{\rm enr}
 \longrightarrow\mathcal C_{\rm cont}^{\rm disc}
\]
est définie et démontrée par les équations, les théorèmes et les reçus imprimés
dans le volume lui-même. Les programmes reproductibles annexés évaluent à nouveau
ces identités ; ils n'apportent ni prémisse absente ni valeur cible.

Les réfutateurs primaires sont : l'échec de l'identité des fibres de
\eqref{eq:frontal-censo-semillas} ; la perte de l'un des champs de
\(C_y\) sous \(U_y\) ; l'incompatibilité entre extension et troncature ;
la vacuité de la limite inverse ; l'absence de contraction ou de surjectivité interne
des cylindres rationnels ; la rupture de naturalité dans la branche
exceptionnelle ; et, pour une publication terminale, le non-respect de
l'hypothèse qui définit son domaine exact. Par conséquent, la validité d'une
sortie se décide dans son domaine et par son réfutateur explicite, non par une
analogie, une nomenclature ultérieure ou une évaluation métrologique.
\end{proposition}

\paragraph{Unicité du système générateur et typage de ses projections.}
Il existe une seule APP, suivie de TRIT et du TPK avec état enrichi.
APP produit les deux feuillets et leurs registres ; TRIT les oriente ; le TPK compose
les voies et accumule feuillet, opposition, résidu, quotient, retenue, bord et
mémoire dans \(\mathcal X_{\rm TPK}^{\rm enr}\). La carte signée est une
projection de ce même état terminal :
\[
 R_{\rm sgn}:\mathcal X_{\rm TPK}^{\rm term,enr}
 \longrightarrow\mathcal U_{12}^{\rm sgn},\qquad
 U_{12}^{\rm sgn}=R_{\rm sgn}(x_{\rm term}),\qquad
 K=\frac14\mathcal H_{12}U_{12}^{\rm sgn},\qquad
 \mathcal H_{12}^2=4I_{12}.
\]
Les deux derniers morphismes sont entiers et réversibles sur leur image. On ne
postule pas d'émetteur APP direct vers \(U_{12}^{\rm sgn}\), car une telle flèche
sauterait TRIT et TPK : le domaine correct de \(R_{\rm sgn}\) est l'état
enrichi déjà généré. Les ombres qui oublient le registre généalogique ou la mémoire sont
des projections ultérieures, non des profils mathématiques alternatifs. Cette unique
composition exclut de ses entrées
\(U,K,D_3K,D_4K,Q,\alpha\) et tout développement cible.

\paragraph{Nombre et particule.}
Un nombre HMT est une section compatible
\(
 \boldsymbol x=(x_n)_{n\geq1}\in
 \varprojlim_n X_n^{\mathrm{enr}}
\)
avec feuillet, orientation, résidu, quotient, retenue, voie, bord et mémoire.
Sa valeur archimédienne est l'unique intersection
\(
 \operatorname{val}(\boldsymbol x)=\bigcap_n I_n(x_n)
\) ;
le quotient par égalité de valeur publie le réel et sa fibre conserve les
généalogies qui le réalisent. Une particule HMT--MD est une classe de sections
persistantes de l'extension enrichie : l'état conserve l'identité,
la voie enregistre son histoire observable et les lecteurs ultérieurs publient
masse, charge, spin et statistique dans leurs codomaines physiques.

\paragraph{Typage du corridor exceptionnel.}
Les lectures de Hadamard et de Paley partent du même état terminal, mais ne sont pas
mises en série l'une avec l'autre. Avec
\(X_\infty^{\rm enr}:=\varprojlim_nX_n^{\rm enr}\), on a des applications distinctes
\begin{equation}
 \operatorname{Pub}_\alpha:X_\infty^{\rm enr}\longrightarrow\mathbb R_{>0},
 \quad x_\infty\longmapsto\alpha_{\rm HMT},
 \qquad
 \operatorname{Inc}_\alpha:X_\infty^{\rm enr}\longrightarrow A_2^{12},
 \quad x_\infty\longmapsto
 v_\alpha=4\rho_1+\rho_2+\cdots+\rho_{12}.
 \label{eq:frontal-dos-publicaciones-alpha}
\end{equation}
Elles partagent une généalogie, mais aucune égalité entre leurs sorties n'est postulée. La branche
exceptionnelle bifurque après \(C_W\) :
\begin{equation}
 \begin{aligned}
 C_P&\longrightarrow A_W=C_P\bmod3\longrightarrow
 C_W=[12,6,6]_3\\
 &\longrightarrow
 \begin{cases}
 W_{12}\xrightarrow{\ \widehat\iota_{D,\ell}\ }
 W_{24}\subset\mathcal G_{24},\\[2pt]
 N(C_W)\xrightarrow{\ \operatorname{Nbr}_{v_\alpha}\ }
 N_\alpha\cong\Lambda_{24}\longrightarrow
 V_{\Lambda_{24}}\longrightarrow V^\natural .
 \end{cases}
 \end{aligned}
 \label{eq:frontal-cadena-excepcional}
\end{equation}
Ici
\begin{equation}
 \mathcal O_{24}:=\{O\subset[24]:|O|=8,\ \mathbf1_O\in\mathcal G_{24}\},
 \qquad
 \mathcal H(D):=\{O\cap D:O\in\mathcal O_{24},\ |O\cap D|=6\}.
 \label{eq:frontal-octadas-hexadas-dodecada}
\end{equation}
\(\ell:W_{12}\simeq(D,\mathcal H(D))\) étant fixé,
\(\widehat\iota_{D,\ell}(H)=O_H\), où
\(O_H\cap D=\ell(H)\). Une telle octade est unique : deux choix contiendraient le
même sous-ensemble de cinq points, en contradiction avec \(S(5,8,24)\). La branche
ternaire recolle \(N(A_2^{12})=N(C_W)\) ; le voisin d'indice trois déterminé
par \(v_\alpha\) produit Leech. Enfin,
\begin{equation}
 \operatorname{Aut}(V^\natural)\cong\mathbb M,\qquad
 a:\mathbb M\times V^\natural\longrightarrow V^\natural,\qquad
 T_g(\tau)=\operatorname{Tr}_{V^\natural}(gq^{L_0-1}).
 \label{eq:frontal-accion-monstruo}
\end{equation}
Les identités \(v_\alpha^2=54\) et \([q]J=196884\) sont des évaluations
scalaires ; aucune ne définit l'action \(a\).

\paragraph{Réalisation parallèle de la dualité T.}
La dualité T ne prolonge pas la branche Moonshine : elle agit dans son propre domaine
\begin{equation}
 \mathscr D_{\rm dual}=\mathbb Z^2\times\mathbb R_{>0},\qquad
 \mathsf T(m,w,R_0)=(w,m,R_0^{-1}),\qquad
 \mathsf T^2=\operatorname{id},\qquad
 E_{R_0}(m,w)=E_{R_0^{-1}}(w,m).
 \label{eq:frontal-dualidad-t}
\end{equation}

\begin{proof}
Le recensement de \eqref{eq:frontal-censo-semillas} construit le catalogue fini et
sélectionne, sans valeurs cibles, les trois blocs initiaux \(w_6\). Le
registre orienté des transitions de l'état TPK complet explicite les douze
blocs postérieurs au germe des trois biographies publiées. Pour chacun des deux
intervalles de transition, les six lignes indépendantes extraites de ces
trois biographies forment des matrices
\(X_\ell,Y_\ell\in M_6(\mathbb F_3)\),
\(\ell\in\{0,1\}\), dont les lignes sont, respectivement, les états de départ
et d'arrivée du registre. Le calcul exact sur ce registre prouve
\(\det X_0=1\), \(\det X_1=-1\) et
\(X_\ell L_\ell=Y_\ell\) ; il reconstruit donc de manière univoque
\(L_\ell=X_\ell^{-1}Y_\ell\). Le
calendrier \((L_0,L_0,L_1,L_1)\) est le seul des seize calendriers
binaires qui conserve simultanément les trois biographies sélectionnées par
les invariants de fibre ; les cent quarante-quatre perturbations
unitaires détruisent cette triple reconnaissance. Par conséquent, les transports,
les douze blocs et leur calendrier sont déduits dans l'unique composition
APP--TRIT--TPK : APP produit les feuillets et le registre de départ, TRIT fixe
l'orientation de chaque transition et TPK prolonge le registre avec mémoire. Le
recensement des germes et l'état enrichi sont des étapes successives du même
générateur, non des profils rivaux ni des théories différentes.

L'extension non scindée
\[
0\longrightarrow C_{12}\longrightarrow C_{108}
\longrightarrow C_9\longrightarrow0
\]
sépare retour de phase et retour d'état. Les troncatures éliminent la
dernière mise à jour et conservent le préfixe ; le lemme de König produit une
branche infinie compatible. Pour chaque précision \(N\), la récurrence et
l'holonomie engendrent la section et ses caractères, et déterminent le cylindre qui
prolonge le registre généalogique précédent. La famille de préfixes satisfait
\[
\rho_{N+1,N}(p_{N+1})=p_N,
\qquad
p_\infty\in\varprojlim_N\{0,\ldots,9\}^{N}.
\]
Mille, dix mille ou un million de chiffres sont des projections finies de cet unique
objet infini.

Les encadrements rationnels emboîtés
\(
 p_N^-<\chi<p_N^+
\)
sont calculés après avoir produit le caractère \(\chi\) et satisfont
\(
 [p_{N+1}^-,p_{N+1}^+]\subset[p_N^-,p_N^+]
\)
avec un diamètre tendant vers zéro. Ils certifient la sortie générée à toute
précision finie ; ils ne sélectionnent ni l'orbite, ni le transport, ni le caractère.

Les trois tours enrichis \(\gamma_{-},\gamma_0,\gamma_+\) de neuf
émissions se distinguent par la retenue
\(\operatorname{Carr}(\gamma_\rho)=\rho\kappa_{\rm Carr}\) et admettent les
trois prolongements après tout préfixe. La limite de leurs voies
enregistrées satisfait
\[
 \operatorname{dig}_{\rm Carr}^{\omega}:
 \Omega_\Gamma^{\rm cal}\xrightarrow{\sim}
 \{-1,0,+1\}^{\mathbb N}.
\]
Après regroupement de six retours, l'application
\[
 \beta=\operatorname{blk}_6\circ
 (\rho\mapsto\rho+1)^{\mathbb N}\circ
 \operatorname{dig}_{\rm Carr}^{\omega}:
 \Omega_\Gamma^{\rm cal}\xrightarrow{\sim}
 (\mathbb F_3^6)^{\mathbb N}
\]
est un homéomorphisme ; son inverse concatène les tours prescrits. Les formes
équilibrées normalisées par
\[
 a_k+c_k=\rho_k+3c_{k+1},\qquad
 \rho_k\in\{-1,0,+1\},
\]
produisent d'abord l'anneau ordonné \(\mathbb Z_{\rm HMT}\), puis son
corps des fractions \(\mathbb Q_{\rm HMT}\) et enfin
\[
 \mathbb R_{\rm HMT}:=
 \operatorname{Cau}(\mathbb Q_{\rm HMT})/{\asymp}.
\]
Si \(\beta(\xi)=(b_q)\) et
\(\nu(b)=\sum_{i=1}^{6}b_i3^{6-i}\), la suite
\[
 q_n=m+\sum_{q=0}^{n-1}\frac{\nu(b_q)}{729^{q+1}}
\]
est de Cauchy, et la partition interne récurrente en \(729\) sous-intervalles
prouve la surjectivité
\[
 P_{\rm ar}:\mathbb Z_{\rm HMT}\times\Omega_\Gamma^{\rm cal}
 \twoheadrightarrow\mathbb R_{\rm HMT},\qquad
 (\mathbb Z_{\rm HMT}\times\Omega_\Gamma^{\rm cal})/{\sim_{\rm ar}}
 \cong\mathbb R_{\rm HMT}.
\]
Ce n'est qu'ensuite que l'unicité du corps ordonné complet archimédien publie
\(\mathbb R_{\rm HMT}\cong\mathbb R\). La direction intérieure subdivise
\([0_{\rm HMT},1_{\rm HMT})\) ; la direction extérieure concatène le préfixe qui code
\(m\in\mathbb Z_{\rm HMT}\). On couvre ainsi la droite sans utiliser au préalable
\(\mathbb R\), de partie entière ni de développement décimal cible.
Les sous-espaces de survie de caractères particuliers individualisent
des sections au sein de cette capacité totale. La couverture du continu provient
de la complétion globale ; la récurrence nonadique engendre
\(\pi_{\mathrm{HMT}},e_{\mathrm{HMT}},\varphi_{\mathrm{HMT}},
\alpha_{\mathrm{HMT}}\), et les filtres typés certifient leurs biographies
au sein de cette couverture.

Les cinq actions spectrale, solénoïdale, de jauge, cohomologique et
déterminantale sont calculées sur une même émission enrichie et partagent
support, voie, registre généalogique et troncatures. Leur naturalité conjointe les conserve
à la limite. La dichotomie parfaite classe chaque publication analytique
comme dénombrable ou de cardinal continu. Le code généré \(h\) sélectionne
l'univers généalogique canonique \(M_h=L[h]\). La construction des
réels HMT identifie chaque nombre par valeur, cylindre, voie et mémoire ; le
bon ordre induit par cette généalogie prouve la formule classique
\[
M_h\models2^{\aleph_0}=\aleph_1.
\]
Comme \(M_h\) est précisément le codomaine engendré par la construction du
continu de ce théorème, la conclusion terminale du système est
\[
 \boxed{2^{\aleph_0}=\aleph_1.}
\]
L'indépendance dans le langage nu de ZFC classe la perte du
sélecteur généalogique \(h\) ; HMT conserve cette donnée et détermine le cardinal de
son continu. L'énoncé conservé est l'hypothèse classique du continu,
avec les cardinaux ordinaires de \(M_h\), et la preuve détaillée établit le
pont entre cylindres compatibles, bon ordre généalogique et cardinalité.
Cette fermeture cardinale se distingue du corridor logique fort, également
construit à partir du continu généalogique :
\[
\begin{aligned}
 \mathcal C_{\mathrm{cont}}^{\mathrm{disc}}
 &\longrightarrow\mathcal L_0^{\mathrm{log}}
 \xrightarrow{\Phi}\mathcal L_1^{\mathrm{log}}
 \longrightarrow\gamma_F\longrightarrow w(\gamma_F)\\
 &\longrightarrow\mathsf{NoEscape}_{24}
 \longrightarrow\mathsf{AntiDiag}_{369}
 \longrightarrow\mathsf{G\ddot{o}delIII}_{\mathrm{HMT}}\\
 &\longrightarrow\mathsf{CycleCriterion}_{\mathrm{ZFC}},
\end{aligned}
\]
qui publie, dans son domaine typé, le critère
\[
 \boxed{
 \neg\operatorname{Consis}(\mathrm{ZFC})
 \Longleftrightarrow
 \exists C\;\bigl(C\text{ est un cycle hors bande et }|C|>24\bigr).}
\]
Les lettres \(\mathcal L_0^{\mathrm{log}},\mathcal L_1^{\mathrm{log}}\)
désignent ici les objets du corridor logique et non les matrices de relèvement
ternaire \(L_0,L_1\) ; la séparation des types empêche qu'une conclusion
locale sur l'un de ces couples affaiblisse celle de l'autre.

Les trois orbites fonctionnelles du catalogue produisent les caractères HMT de
clôture, de propagation et d'auto-échelle. Leurs lois internes imposent,
respectivement,
\[
16\arctan(1/5)-4\arctan(1/239)=\pi_{\mathrm{HMT}}=\pi,
\qquad
P_{s+t}=P_sP_t\Rightarrow
P_1=\sum_{n\ge0}\frac1{n!}=e_{\mathrm{HMT}}=e,
\]
\[
F_{\mathrm{av}}\binom1{\varphi_{\mathrm{HMT}}}
=\varphi_{\mathrm{HMT}}\binom1{\varphi_{\mathrm{HMT}}},
\qquad \varphi_{\mathrm{HMT}}=\varphi.
\]
L'involution interne échange les branches
\(\pi_{\mathrm{HMT}}/e_{\mathrm{HMT}}\), conserve l'axe
\(\varphi_{\mathrm{HMT}}\) et retient l'agrégat
\(e_{\mathrm{HMT}}+\pi_{\mathrm{HMT}}\). Dans la phase TPK enrichie de
l'unique composition, la projection canonique de l'état terminal publie sa coordonnée primitive
signée
\(U=R_{\mathrm{sgn}}(x_\infty)\), et l'identité
\(\mathcal H_{12}^2=4I\) reconstruit de manière entière et unique le sceau
\(K=\mathcal H_{12}U/4\). La récurrence dodécaphasique avec retenue et le
prolongement distant de ses cylindres publient \(\alpha_{\rm HMT}\). De
manière indépendante, l'état enrichi détermine la signature
\[
 \mathcal I_9=(2\pi_{\rm HMT},3,4,7,20,21,45,46,120,11),
\]
et le lecteur gradué associé construit, sans consulter la racine, l'opérateur
polynomial exact
\(E_{21/22}^{(9)}=P_9-\log_{10}(10/9)\). Son unique racine positive simple
\(\xi_9\) est le témoin fini d'ordre neuf de la seconde publication
analytique interne. Le même transport gradué du TPK est itéré sur les
degrés successifs et construit le système compatible complet
\(E_{21/22}^{(6+3m)}\) ; ses carrés de troncature commutent avec ceux de la
famille dodécaphasique et déterminent, à chaque profondeur, le même cylindre
survivant. En conséquence,
\[
 \alpha_A:=\varprojlim_N\rho_N(\alpha_{12}),\qquad
 E_{21/22}:=\varprojlim_mE_{21/22}^{(6+3m)},\qquad
 \alpha_B:=\operatorname*{arg\,zero}_{x\in I_\alpha}E_{21/22}(x),
\]
\[
 \boxed{\alpha_A=\alpha_B=\alpha_{\rm HMT}.}
\]
L'égalité finie
\[
 \operatorname{Tr}_9(\xi_9^{-1})
 =\operatorname{Tr}_9(\alpha_{12}^{-1})=137.035999084
\]
certifie le premier carré visible de cette identité ; elle ne la définit pas, ne la
limite pas à neuf chiffres et ne sélectionne pas les coefficients du prolongement.
La clôture bidirectionnelle fonctionnelle associée est
\[
\alpha T^2-B_9(\alpha)T+\alpha^3=0,
\qquad
T_\pm=\alpha e^{\pm\xi},
\qquad
T_+T_-=\alpha^2,
\qquad
J_\alpha(T)=\frac{\alpha^2}{T}.
\]
Ainsi, la coordonnée publiée par la clôture dodécaphasique est la moyenne
géométrique conservée par les deux voies conjuguées. La voie dodécaphasique et le
lecteur analytique maintiennent leurs domaines, opérateurs et réfutateurs propres, mais
leurs systèmes inverses publient une même section limite. Le comparateur
fini et la reconnaissance métrologique sont ultérieurs.

La coordonnée d'incidence du secteur \(\alpha\) fixe
\[
v_\alpha=4\rho_1+\rho_2+\cdots+\rho_{12},
\qquad v_\alpha^2=54,
\]
et le voisin d'indice trois
\[
N_{\alpha,0}
=\{x\in N(C_W):\langle x,v_\alpha\rangle\equiv0\pmod3\},
\qquad
N_\alpha=N_{\alpha,0}+\mathbb Z\frac{v_\alpha}{3}
\cong\Lambda_{24}.
\]
Cette branche construit exactement le voisin de Leech, puis la chaîne
\(\Lambda_{24}\to V^\natural\) et l'action de
\(\operatorname{Aut}(V^\natural)\cong\mathbb M\) sur \(V^\natural\). Le
scalaire \(\alpha_{\mathrm{HMT}}\) et le vecteur \(v_\alpha\) sont les deux
coordonnées typées du même secteur dodécaphasique. Leur couplage exact est
la charnière d'origine commune qui conduit aux symétries du Monstre : la
coordonnée scalaire est publiée dans \(\mathbb R_{\rm HMT}\) et la coordonnée
d'incidence est publiée dans la branche
\(A_2^{12}\to\Lambda_{24}\to V^\natural\to\mathbb M\), avec leurs codomaines
conservés et leur antécédent généalogique partagé explicite.

La mécanique dimensionnelle reçoit les sorties précédentes et compose, dans leur ordre
constitutif, la rétention de phase, la décade \(-34\), le couple de Planck,
les canaux \(90/120\), l'électron, l'atlas des masses, la gravité, l'aire
et l'information. Les chapitres suivants démontrent chaque flèche, publient
son réfutateur et relient son certificat exécutable. Les identifications
conventionnelles apparaissent ensuite comme reconnaissance des grandeurs déjà
produites. Cela complète la chaîne de
\eqref{eq:frontal-doble-proyeccion}.
\end{proof}

\section*{Hiérarchie des publications géométriques et dimensionnelles}
\addcontentsline{toc}{section}{Hiérarchie des publications géométriques et dimensionnelles}

\begin{proposition}[TRIT comme unité locale des trois régimes]
La signature tritique conserve deux lecteurs coordonnés. Le lecteur géométrique
publie \(\tau_{\rm g}\in\{+1,0,-1\}\) et
\begin{equation}
 J_{\tau_{\rm g}}^{\,2}=-\tau_{\rm g}I,\qquad
 J_+^{\,2}=-I,\qquad J_0^{\,2}=0,\qquad J_-^{\,2}=+I,
 \label{eq:frontal-trit-tres-regimenes}
\end{equation}
et distingue les régimes elliptique, parabolique et hyperbolique. Le lecteur
temporel actif publie un autre signe, \(\varepsilon_{\rm t}\), au moyen de
\begin{equation}
 \begin{aligned}
 +1&\mapsto\text{retour, mémoire et passé},\\
 0&\mapsto\text{seuil et présent},\\
 -1&\mapsto\text{ouverture bidirectionnelle et futur}.
 \end{aligned}
 \label{eq:frontal-trit-tiempo}
\end{equation}
On ne postule pas \(\tau_{\rm g}=\varepsilon_{\rm t}\). Ces deux sorties
n'identifient pas signe de courbure et orientation temporelle : ce sont des projections
distinctes d'une même unité locale. Le TPK transporte
l'unité complète dans
\begin{equation}
 \mathcal F_q=
 \mathcal O_q\times\mathcal H_q\times\mathcal C_q\times
 \mathcal S_q\times\mathcal B_q\times\Lambda_q
 \label{eq:frontal-fibra-tpk}
\end{equation}
et conserve feuillet et orientation, résidu et quotient, retenue, signature, cylindre,
bord et registre. En particulier, la branche elliptique possède le réalisateur
fini
\begin{equation}
 A_W^2=2I_6=-I_6\quad\text{dans }M_6(\mathbb F_3),\qquad
 i_{\rm HMT}:=\chi_i(A_W),\qquad i_{\rm HMT}^{\,2}=-1,
 \label{eq:frontal-raiz-interna-menos-uno}
\end{equation}
où \(\chi_i:\langle A_W\rangle\to\mathbb C\) est la carte orientée.
La racine de \(-1\), le plan complexe et le système d'Euler sont des publications
d'un endomorphisme construit, non des entrées du générateur.
\end{proposition}

\begin{proof}
Le sélecteur tritique incorpore l'un des trois générateurs de
\eqref{eq:frontal-trit-tres-regimenes} sans effacer les feuillets somme et produit.
Le transport cardinal conserve son orientation et la mise à jour conserve
les six facteurs de \eqref{eq:frontal-fibra-tpk}. Par conséquent, la projection
réelle peut condenser les trois régimes, tandis que l'état enrichi
reconstruit lequel d'entre eux et quelle mémoire ont produit la valeur visible. Les
réalisations trigonométrique, nilpotente et hyperbolique développées dans le
corps du texte sont les images respectives de cet unique typage local.
\end{proof}

\begin{theorem}[Chaîne d'action, de gravité et de matière]
Après publication corrélée de
\((\pi_{\rm HMT},\varphi_{\rm HMT},e_{\rm HMT},\alpha_{\rm HMT})\), la branche
déterminantale fixe la décade
\begin{equation}
 \Sigma_H=\varphi_{\rm HMT}\alpha_{\rm HMT}^{16},\qquad
 d_H:=\lfloor\log_{10}\Sigma_H\rfloor=-34.
 \label{eq:frontal-decada-accion}
\end{equation}
La mantisse provient du lecteur régional distinct
\begin{align}
 H_5&=\frac12\sqrt{\frac{1000\alpha_{\rm HMT}}{\varphi_{\rm HMT}}}
 -\alpha_{\rm HMT}+\frac9{16}\alpha_{\rm HMT}^2
 -\frac59\alpha_{\rm HMT}^3+\frac7{48}\alpha_{\rm HMT}^4
 -\frac1{54}\alpha_{\rm HMT}^5,\notag\\
 D_A&=\exp\!\left(-\frac{100\pi_{\rm HMT}}9\alpha_{\rm HMT}\right),
 \qquad
 \mathscr C_\pi=169\alpha_{\rm HMT}^6+
 \frac{D_A\alpha_{\rm HMT}^7}{1-D_A\alpha_{\rm HMT}},\notag\\
 \eta_{\rm pre}&:=H_5,\qquad
 \eta_{\rm ret}:=H_5-\frac{90}{\pi_{\rm HMT}}\mathscr C_\pi>0.
 \label{eq:frontal-mantisas-accion}
\end{align}
Pour \(a\in\{\mathrm{pre},\mathrm{ret}\}\), les deux sections d'action sont
\begin{equation}
 \hbar_a=\eta_a10^{d_H}\mathcal U_S
 =\eta_a10^{-34}\mathcal U_S,\qquad
 h_a=2\pi_{\rm HMT}\hbar_a.
 \label{eq:frontal-secciones-accion}
\end{equation}
\(\Sigma_H\) fixe seulement la décade ; il n'est pas présenté comme générateur de la
mantisse.

\(t_0>0\), \(c_{\rm int}>0\) et \(\ell_0=c_{\rm int}t_0\) étant fixés, CT108
publie
\begin{equation}
 \omega_{108}=\frac{2\pi_{\rm HMT}}{108t_0},\qquad
 R_{108}=\frac{c_{\rm int}}{\omega_{108}}
 =\frac{54}{\pi_{\rm HMT}}\ell_0,\qquad
 m_{108,a}=\frac{\hbar_a\omega_{108}}{c_{\rm int}^2}.
 \label{eq:frontal-reloj-gravedad}
\end{equation}
La famille gravitationnelle typée est
\begin{equation}
 G_a(\vartheta)=\vartheta\frac{c_{\rm int}^5t_0^2}{\hbar_a}.
 \label{eq:frontal-familia-gravedad}
\end{equation}
Dans la convention réduite
\(r_{g,a}=G_a m_{108,a}/c_{\rm int}^2\), la fermeture
\(r_{g,a}(\vartheta)=R_{108}\) possède l'unique solution positive
\begin{equation}
 \vartheta_{108}=\left(\frac{54}{\pi_{\rm HMT}}\right)^2,\qquad
 G_a:=G_a(\vartheta_{108}),
 \label{eq:frontal-selector-gravedad}
\end{equation}
et alors seulement publie
\begin{equation}
 \ell_P:=\sqrt{\frac{G_a\hbar_a}{c_{\rm int}^{\,3}}}
 =\sqrt{\vartheta_{108}}\,\ell_0
 =R_{108}=\bar\lambda_{C,a}=r_{g,a}.
 \label{eq:frontal-longitud-planck}
\end{equation}
En outre, \(G_a\hbar_a=\vartheta_{108}c_{\rm int}^5t_0^2\), de sorte que
\(\ell_P\) est indépendant de l'orientation \(a\). La gravité tensorielle
\[
 \mathbb G_{{\rm TT},a}(n)
 =G_a\Lambda(n)\Gamma_5P_5:
 V_{12}\longrightarrow\mathcal H_{\rm TT}(n)
\]
est postérieure au scalaire. L'action Palatini--Cartan--Holst reçoit ensuite
\(G_a\) et la sortie indépendante \(\gamma_{\rm HMT}\) ; la fonctionnelle
hexade--octade transporte cette dernière vers le spectre d'aire et l'information
de bord.

Dans la branche de matière, la loi complète agit sur \(81\) cellules, \(56\)
familles, \(13\) multisections et \(324\) voies, avec un secteur nul séparé. Les
lecteurs \(K_D\) et \(K_\Omega\) possèdent, respectivement, \(14\) et \(9\)
profils, forment \(40\) couples et coïncident sur \(18\) voies uniquement dans
\((80,54,6)\). La fibre centrale \((5,5)\) produit l'électron bêta. Sur les
fibres restantes, la loi publie des opérateurs positifs, des mesures spectrales et des
multiplicités ; la reconnaissance extérieure s'ouvre ensuite au moyen de
l'inventaire exhaustif
\begin{equation}
 \mathscr I_{\rm ext}
 =\mathscr I_m^{471}\coprod\mathscr I_\Gamma^{384},
 \qquad |\mathscr I_{\rm ext}|=855.
 \label{eq:frontal-inventario-masas-anchuras}
\end{equation}
CKM transporte les repères spectraux quark et PMNS les repères chargé et neutre
de rang trois ; leurs domaines, opérateurs et réfutateurs demeurent distincts.
Aucune valeur de l'inventaire ne sélectionne une voie, un coefficient, une valeur propre
ou un angle.
\end{theorem}

\begin{proof}
La forme normale de Smith produit
\(\operatorname{diag}(1,2,2,4)\) et un conoyau d'ordre \(16\) ; les bornes de
\(\Sigma_H\) fixent \(d_H=-34\), tandis que
\eqref{eq:frontal-mantisas-accion}--\eqref{eq:frontal-secciones-accion}
publient \(\hbar_a\) puis \(h_a\). La substitution de
\(\vartheta_{108}\) dans \(G_a(\vartheta)\) prouve
\eqref{eq:frontal-longitud-planck}. L'incidence des ensembles
\(H_{\rm act}\times\binom{H_{\rm act}}2\) et
\(\{1,\ldots,8\}\times\binom{H_{\rm act}}2\) produit \(90\) et \(120\).
Les douze secteurs dodécaphasiques et l'unique couture orientée de retour
produisent, au moyen du lecteur linéaire, le poids \(12{:}{-}1\) de la transition
hexade--octade. Le coefficient de \((1-q)^{-1}\), égal à \(1/24\), et la
minimalité diophantienne vérifient ensuite ce poids. Le recensement fini de l'atlas et de ses deux lecteurs
prouve les cardinaux de la branche de masse ; l'inversion cellulaire possède un unique
point fixe, \((5,5)\), qui réduit canoniquement la fibre électronique au rang
un. L'inventaire externe n'est composé qu'après avoir figé la fibre,
l'opérateur, le projecteur, la représentation, le schéma et l'échelle, de sorte qu'une
permutation de ses grandeurs laisse invariant tout le constructeur interne.
\end{proof}
% END HMT-MD-FRONTAL-MAXIMAL-SIGNATURES-20260904
